\documentclass[10pt,a4paper]{amsart}
\usepackage[margin=19mm]{geometry}
\usepackage{amsmath,amssymb,mathtools}
\usepackage[scr=rsfs,cal=euler]{mathalfa}
\usepackage{xfrac}
\usepackage[unicode,hidelinks,bookmarks=false]{hyperref}
\newcommand{\ie}{i.e.\ }
\newcommand{\cf}{cf.\ }
\newcommand{\resp}{respectively\ }
\newcommand{\Subsection}{\subsection}
\providecommand{\nobreakdash}{\nobreak\hbox{-}}
\setlength{\emergencystretch}{2em}
\multlinegap0pt
\usepackage{amssymb}
\usepackage{mathtools}
\usepackage{enumerate}
\usepackage{xcolor}
\newtheorem{lemma}{Lemma}
\newtheorem{theorem}{Theorem}
\newcommand{\Prob}{\mathbb{P}}
\title{A Heuristic approach to the Iwasawa theory of elliptic curves}
\author{Katharina M\"uller, Anwesh Ray}
\date{Source: arXiv:2409.15056v2 (CC BY 4.0)}
\begin{document}
\maketitle

\noindent\emph{Source and licence.} A Heuristic approach to the Iwasawa theory of elliptic curves. Version arXiv:2409.15056v2 (CC BY 4.0), distributed by arXiv under the Creative Commons Attribution 4.0 International licence, http://creativecommons.org/licenses/by/4.0/, as recorded in the arXiv licence metadata for this exact version and shown on the abstract page https://arxiv.org/abs/2409.15056v2. Full licence notice: \texttt{licenses/arxiv-2409.15056-CC-BY-4.0.txt}.\par\smallskip

\noindent\textit{Editorial scope.} Counted unit: the article's theorem thm:probability (source0.tex lines 628--634). The article's complete original proof of it is delivered with it (source0.tex lines 635--653, one \texttt{proof} environment of 19 lines). No mathematics is written, repaired or re-derived: the statement and the proof are the article's own lines, in the article's own notation, and no retained line is substituted or re-flowed.  Retained uncounted as the source-local context the proof needs: lines 583--586; lines 597--604.  Nothing was omitted from any retained span, so the package carries no omission record.  No retained line contains a citation, so the wrapper carries no bibliography and no bibliography entry.  No retained line contains a figure environment, an image-inclusion command or drawing code, so no image is delivered and the package declares no asset dependency.  Editorial formatting only, in addition: an AMS \texttt{amsart} preamble that restates the article's own declarations and macros used by the retained lines, namely the environments \texttt{lemma}, \texttt{theorem} on separate counters, together with the standard packages those lines need.  The retained environments print in this document as Lemma 1, Theorem 1 in order of appearance, whereas the article prints the same items with the numbering of its own full document.  The complete native source of the article as distributed in the arXiv e-print for this exact version is bundled unchanged as \texttt{sources/165677/source0.tex}.
\begin{lemma}
\label{number of maximal submodules}
    With respect to notation above, there are $p^{n-1}(p+1)$ maximal submodules in $\Omega_n^2$.
\end{lemma}

\begin{lemma}
    \label{max-cyclic}
    Let $N_1$ and $N_2$ be maximal submodules. Then one of the following is true
    \begin{itemize}
        \item $N_1\cap N_2=0$.
        \item $N_1=N_2$.
    \end{itemize}
\end{lemma}

\begin{theorem}
\label{thm:probability}
    Let $N_1,N_2$ be cyclic submodules of $\Omega^2$. Then we have
    \[\Prob(\{(N_1, N_2)\in \mathcal{M}\mid N_1\cap N_2=0\})=1.\]
    %Furthermore,
   % \[\sup_{\substack{B\subset \{(N_1, N_2)\in \mathcal{M}\mid N_1\cap N_2=0\}\\B\in \mathcal{A}}}\Prob(B)=1.\]
\end{theorem}

\begin{proof}
By Lemma \ref{max-cyclic} we obtain 
    \begin{align*}
        \Prob(N_1\cap N_2=0)=\Prob(N_1\neq N_2)\\
        \ge \Prob(\pi_{n}(N_1)\neq \pi_n(N_2)).\\
    \end{align*}
    Note that the condition $\pi_n(N_1)\neq \pi_n(N_2)$ actually produces a measurable set. Therefore, we find that 
    \[\begin{split} & \Prob(\pi_{n}(N_1)\neq \pi_n(N_2)) \\
    = & \Prob_n(\pi_{n}(N_1)\neq \pi_n(N_2)) \\
    = &1-\Prob_n(\pi_{n}(N_1)= \pi_n(N_2)).\end{split}\]
    Let $\mathcal{N}_n$ be the set of maximal submodules in $\Omega^2$. 
     By Lemma \ref{number of maximal submodules} we have
    \begin{align*}\Prob_n(\pi_{n}(N_1)=\pi_n(N_2))=\sum_{M\in \mathcal{N}_n}\Prob_n(\pi_n(N_1)=M)\Prob(\pi_n(N_2)=M)\\=\sum_{M\in \mathcal{N}_n}\left(\frac{1}{(p+1)p^{n-1}}\right)^2=\frac{1}{p^{n-1}(p+1)} .\end{align*}
    This implies that 
    \[\Prob(N_1\cap N_2=0)\ge 1-\frac{1}{(p+1)p^{n-1}}\]
    for all $n$. The left hand side of the above inequality does not depend on $n$, thus, letting $n$ tend to infinity gives the claim.

    %We have already seen that  \[\{(N_1, N_2)\in \mathcal{M}\mid N_1\cap N_2=0\}\supset \{(N_1,N_2)\mid \pi_n(N_1)\neq \pi_n(N_2)).\] The right hand side is clearly measurable and has measure $1-\frac{1}{(p+1)p^{n-1}}$ (as we computed above). Taking the suprenum over $n$ proves the second claim.
\end{proof}

\end{document}
